\section{More stabilizer formalism}\label{sec:more-stabilizer-formalism}

If we were sticking with static stabilizer codes, we'd have already said everything we need to about the stabilizer formalism.
But to be able to talk about dynamic stabilizer codes, we need to know a bit more about how $\VVV_\SSS$ and $\mathcal{N(S)/S}$ evolve under applying Cliffords and measuring Paulis.

\subsection{Evolution of $\VVV_\SSS$}\label{subsec:evolution-of-codespace}

The result of applying a Clifford or measuring a Pauli and getting outcome $m$ is that a linear map $A_m$ is applied to our system of $n$ qubits -- i.e.\ to the Hilbert space $\mathcal{H}$.
If a Clifford unitary $U$ is applied, the map $A_m$ is simply this Clifford.
If a Pauli $P$ is measured, and outcome $m$ obtained, then the map $A_m$ is the projector $\Pi_{mP}$.
In this case, after $A_m$ is applied, the post-measurement state is also renormalised.
We are only interested in the restriction $A_m|_{\VVV_\SSS} : \VVV_\SSS \to \VVV_{\SSS'}$ of this map to the codespace $\VVV_\SSS$.

\teague{Add figure of all four cases}

The most straightforward case is when $A_m$ is a Clifford unitary $U$.
Then the restriction $A_m|_{\VVV_\SSS}$ is a unitary map $\VVV_\SSS \to \VVV_{\SSS'}$.
Not too much more complicated is \hyperref[thm:stabilizer_group_update_measurement:case1]{Case 1} of the Pauli measurements, where $\pm P \in \mathcal{S}$ and $A_m = \Pi_{\pm P}$.
Given a state $\ket{\psi} \in \SSS$, we have $A_m \ket{\psi} = \Pi_{\pm P} \ket{\psi} = \ket{\psi}$, hence $A_m|_{\VVV_\SSS}$ is in fact the identity when restricted to this subspace.
In \hyperref[thm:stabilizer_group_update_measurement:case2]{Case 2} we measure a $P$ that anti-commutes with exactly one generator $S_1$ of $\SSS$, with random outcome $m$ and map $A_m = \Pi_{mP}$.
Since the outcome $m$ is uniformly random across all states in $\VVV_\SSS$, the post-measurement renormalisation is the same for all states -- we just multiply by $1/{\sqrt{1/2}} = \sqrt{2}$.
Hence the overall map applied to $\VVV_\SSS$ is $\sqrt{2} A_m|_{\VVV_\SSS}$
We can show that this is in fact a unitary $\VVV_\SSS \to \VVV_{\SSS'}$.

\begin{proposition}\label{prop:A_m_unitary_up_to_normalisation}
    In \hyperref[thm:stabilizer_group_update_measurement:case2]{Case 2}, the restriction of $\sqrt{2} A_m$ to $\VVV_\SSS$ is a unitary $\VVV_\SSS \to \VVV_{\SSS'}$.
    \begin{proof}
        Will show $\sqrt{2} A_m|_{\VVV_\SSS}$ really does map states in $\VVV_\SSS$ to states in $\VVV_{\SSS'}$, and also preserves inner products, making it an isometry $\VVV_\SSS \to \VVV_{\SSS'}$.
        But since we know the dimensions of $\VVV_\SSS$ and $\VVV_{\SSS'}$ are equal, this in fact makes $\sqrt{2} A_m|_{\VVV_\SSS}$ a unitary.
        So, to show $im(\sqrt{2} A_m|_{\VVV_\SSS}) \leq \VVV_{\SSS'}$, it suffices to show that for any $\ket{\psi} \in \VVV_\SSS$ its image $\sqrt{2} A_m\ket{\psi}$ is stabilized by all generators of $\SSS' = \langle mP, S_2, \ldots, S_r \rangle$.
        First, show $mP$ stabilizes $\sqrt{2} A_m\ket{\psi}$:
        \begin{equation}\label{eq:mP_stabilizes_new_state}
        \begin{aligned}
            mP(\sqrt{2} A_m\ket{\psi})
            &= mP(\sqrt{2} \cdot \frac{1}{2}(I + mP))\ket{\psi} \\
            &= \sqrt{2} \cdot \frac{1}{2}(mP + m^2 P^2) \ket{\psi} \\
            &= \sqrt{2} \cdot \frac{1}{2}(mP + I) \ket{\psi} \\
            &= \sqrt{2} A_m\ket{\psi} \\
        \end{aligned}
        \end{equation}
        For the remaining stabilizers, since $[P, S_j] = 0$ for all $j \geq 2$, we get:
        \begin{equation}\label{eq:S_j_stabilizes_new_state}
        \begin{aligned}
            S_j(\sqrt{2} A_m\ket{\psi})
            &= S_j(\sqrt{2} \cdot \frac{1}{2}(I + mP))\ket{\psi} \\
            &= (\sqrt{2} \cdot \frac{1}{2}(I + mP)) S_j \ket{\psi} \\
            &= \sqrt{2} \cdot \frac{1}{2}(I + mP) \ket{\psi} \\
            &= \sqrt{2} A_m\ket{\psi} \\
        \end{aligned}
        \end{equation}
        Now just need to show $\sqrt{2} A_m|_{\VVV_\SSS}$ preserves inner products.
        Consider $\ket{\psi}, \ket{\phi} \in \VVV_\SSS$.
        First, note that $\braket{\psi | P | \phi} = 0$ because it's equal to its negation:
        \begin{equation}\label{eq:P_phi_orthogonal_to_psi}
        \braket{\psi | P | \phi}
        = \braket{\psi | S_1 P S_1| \phi}
        = - \braket{\psi | P S_1 S_1| \phi}
        = - \braket{\psi | P | \phi}
        \end{equation}
        Therefore $\bra{\psi} (\sqrt{2} A_m)^\dagger (\sqrt{2} A_m) \ket{\phi} = \braket{\psi | \phi}$:
        \begin{equation}\label{eq:A_preserves_inner_products}
        \begin{aligned}
            \braket{\psi | (\sqrt{2} A_m)^\dagger (\sqrt{2} A_m) | \phi}
            &= \braket{\psi | 2 \Pi_{mP}^\dagger \Pi_{mP} | \phi} \\
            &= \braket{\psi | 2\Pi_{mP} | \phi} \\
            &= \braket{\psi | (I + mP) | \phi} \\
            &= \braket{\psi | I | \phi} + 0\\
            &= \braket{\psi | \phi} \\
        \end{aligned}
        \end{equation}
    \end{proof}
\end{proposition}

Finally, in \hyperref[thm:stabilizer_group_update_measurement:case3]{Case 3}, we measure a non-trivial logical operator $P \in \NNN(\SSS) - (\pm \SSS)$, getting outcome $m$ and map $A_m = \Pi_{mP}$.
This continues to be a projection when restricted just to the codespace $\VVV_\SSS$; its image is exactly the new codespace $\VVV_{\SSS'}$.

\begin{proposition}
    In \hyperref[thm:stabilizer_group_update_measurement:case3]{Case 3}, the restriction of $A_m$ to $\VVV_\SSS$ is a projector $\VVV_\SSS \to \VVV_{\SSS'}$
    \begin{proof}
        Will again show $A_m|_{\VVV_\SSS}$ really does map states in $\VVV_\SSS$ to states in $\VVV_{\SSS'}$, and then show it's self-adjoint and idempotent, i.e.\ a projector.
        Can prove the former exactly as in the previous proposition, by showing that for any $\ket{\psi} \in \VVV_\SSS$ its image $A_m\ket{\psi}$ is stabilized by all generators of $\SSS' = \langle mP, S_1, S_2, \ldots, S_r \rangle$, where $[P, S_j] = 0$ for all $j$.
        Then to prove the latter, first note that we can write $A_m|_{\VVV_\SSS}$ in terms of the codespace projector $\Pi_\SSS \coloneqq \prod_{j} \Pi_{S_j} = \prod_{j} \frac{1}{2}(I + S_j)$.
        Namely, we have $A_m|_{\VVV_\SSS} = A_m \Pi_\SSS$.
        But since the image of this is in $\VVV_{\SSS'} \leq \VVV_\SSS$, we can add an extra $\Pi_\SSS$ term: $A_m|_{\VVV_\SSS} = \Pi_\SSS A_m \Pi_\SSS$.
        Then since $A_m = \Pi_{mP}$ and $\Pi_{\SSS}$ are both projectors, they're self-adjoint.
        So $A_m|_{\VVV_\SSS}$ is too:
        \begin{equation}\label{eq:A_m_restriction_self_adjoint}
            A_m|_{\VVV_\SSS}^\dagger
            = (\Pi_\SSS A_m \Pi_\SSS)^\dagger
            = \Pi_\SSS^\dagger A_m^\dagger \Pi_\SSS^\dagger
            = \Pi_\SSS A_m \Pi_\SSS
            = A_m|_{\VVV_\SSS}
        \end{equation}
        Then show it's idempotent; note that since $[P, S_j] = 0$ for all $j$, the map $A_m = \Pi_{mP} = \frac{1}{2}(I + mP)$ commutes with the codespace projector $\Pi_\SSS = \prod_{j} \Pi_{S_j} = \prod_{j} \frac{1}{2}(I + S_j)$.
        Combining this with idempotence of projectors $A_m$ and $\Pi_{\SSS}$, we get:
        \begin{equation}\label{eq:A_m_restriction_idempotent}
%            \begin{aligned}
%                A_m|_{\VVV_\SSS}^2
%                &= (A_m \Pi_\SSS)(A_m \Pi_\SSS) \\
%                &= A_m A_m \Pi_\SSS \Pi_\SSS \\
%                &= A_m \Pi_\SSS \\
%                &= A_m|_{\VVV_\SSS}
%            \end{aligned}
            A_m|_{\VVV_\SSS}^2
            = (A_m \Pi_\SSS)(A_m \Pi_\SSS)
            = A_m A_m \Pi_\SSS \Pi_\SSS
            = A_m \Pi_\SSS
            = A_m|_{\VVV_\SSS}
        \end{equation}
    \end{proof}
\end{proposition}

We can package all of the above into a single theorem.

\begin{theorem}\label{thm:codespace_maps}
    Given a stabilizer group $\SSS$, the effect of applying a Clifford and \hyperref[thm:stabilizer_group_update_measurement:case1]{Cases 1} and \hyperref[thm:stabilizer_group_update_measurement:case2]{2} of measuring a Pauli on the codespace $\VVV_\SSS$ is that it is mapped unitarily to the new codespace $\VVV_{\SSS'}$.
    In \hyperref[thm:stabilizer_group_update_measurement:case3]{Case 3}, it is mapped projectively onto the new codespace $\VVV_{\SSS'}$.
\end{theorem}

\begin{example}\label{ex:bacon_shor_codespace_updates}
    In \Cref{ex:bacon_shor_stabilizer_group_updates}, we started with a trivial stabilizer group $\SSS_{-1}$ on $4$ qubits and repeatedly measured $X_1 X_2$, $X_3 X_4$, $Z_1 Z_3$ and $Z_2 Z_4$, in that order.
    So initially the codespace is $\HHH$, the whole Hilbert space.
    Both of the first two measurements $X_1 X_2$ and $X_3 X_4$ are in \hyperref[thm:stabilizer_group_update_measurement:case3]{Case 3}, so the resulting maps are projections that each halve the dimension of the codespace:
    \begin{equation}
        \VVV_{\SSS_{-1}}
        \xrightarrow{\Pi_{a_0 X_1 X_2}}
        \VVV_{\SSS_0}
        \xrightarrow{\Pi_{b_1 X_3 X_4}|_{\VVV_{\SSS_{0}}}}
        \VVV_{\SSS_0'}
    \end{equation}
    The next measurement $Z_1 Z_3$ is in \hyperref[thm:stabilizer_group_update_measurement:case2]{Case 2}, so maps us unitarily to the next codespace $\SSS_0''$.
    Then the measurement $Z_2 Z_4$ is in \hyperref[thm:stabilizer_group_update_measurement:case3]{Case 3}, so another projector is applied, mapping from $\SSS_0''$ to the new codespace $\SSS_0'''$, which is again half the size:
    \begin{equation}
        \VVV_{\SSS_{0}'}
        \xrightarrow{\sqrt{2}\Pi_{c_0 Z_1 Z_3}|_{\VVV_{\SSS_{0}'}}}
        \VVV_{\SSS_0''}
        \xrightarrow{\Pi_{d_0 Z_2 Z_4}|_{\VVV_{\SSS_{0}''}}}
        \VVV_{\SSS_0'''}
    \end{equation}
    Now we start the next round of measurements.
    Measuring $X_1 X_2$ is this time in \hyperref[thm:stabilizer_group_update_measurement:case2]{Case 2}.
    Writing the new stabilizer group as $\SSS_1$, this measurement thus maps us unitarily from $\VVV_{\SSS_0'''}$ to $\VVV_{\SSS_1}$.
    Measuring $X_3 X_4$ is in \hyperref[thm:stabilizer_group_update_measurement:case1]{Case 1}, so at the level of the codespace just applies the identity map to $\VVV_{\SSS_1}$.
    The same holds analogously when measuring $Z_1 Z_3$ and $Z_2 Z_4$; the first of these is in \hyperref[thm:stabilizer_group_update_measurement:case2]{Case 2} and maps unitarily from $\VVV_{\SSS_1}$ to $\VVV_{\SSS_1'}$, then the second is in \hyperref[thm:stabilizer_group_update_measurement:case1]{Case 1} and applies the identity to $\VVV_{\SSS_1'}$.
    This pattern repeats for every subsequent round of measurements.
    So the measurements from round 1 onwards implement a sequence of unitary maps:
    \begin{equation}
        \VVV_{\SSS_0'''}
        \xrightarrow{\sqrt{2}\Pi_{a_1 X_1 X_2}|_{\VVV_{\SSS_0'}}}
        \VVV_{\SSS_1}
        \xrightarrow{\sqrt{2}\Pi_{c_1 Z_1 Z_3}|_{\VVV_{\SSS_1}}}
        \VVV_{\SSS_1'}
        \xrightarrow{\sqrt{2}\Pi_{a_1 X_1 X_2}|_{\VVV_{\SSS_1'}}}
        \VVV_{\SSS_2}
        \ldots
    \end{equation}
\end{example}

\begin{example}\label{ex:teleportation_codespace_updates}
    Recall the teleportation example of \Cref{ex:teleportation_stabilizer_group_updates}.
    Starting from $\SSS = \langle X_2 X_3, Z_2 Z_3 \rangle$, we measured $X_1 X_2$ then $Z_2 Z_3$.
    Both of these were \hyperref[thm:stabilizer_group_update_measurement:case2]{Case 2} measurements, so together they unitarily map us from $\VVV_\SSS$ to $\VVV_{\SSS''}$:
    \begin{equation}
        \VVV_{\SSS}
        \xrightarrow{\sqrt{2}\Pi_{a X_1 X_2}|_{\VVV_{\SSS}}}
        \VVV_{\SSS'}
        \xrightarrow{\sqrt{2}\Pi_{b Z_1 Z_2}|_{\VVV_{\SSS'}}}
        \VVV_{\SSS''}
    \end{equation}
    We then applied Pauli correction $U = Z_3^{\frac{1 - a}{2}} X_3^{\frac{1 - b}{2}}$.
    Since this is in particular a Clifford unitary, this too applies a unitary to our codespace:
    \begin{equation}
        \VVV_{\SSS''}
        \xrightarrow{U|_{\VVV_{\SSS''}}}
        \VVV_{\SSS'''}
    \end{equation}
    But as in \Cref{ex:teleportation_stabilizer_group_updates} this is still not enough to conclude that this teleportation scheme works as advertised!
    We want to show that the resulting unitary map is in some sense `the identity map' from $\VVV_{\SSS}$ to $\VVV_{\SSS'''}$.
    That is, if we've picked particular operators to represent logical $X$ and $Z$ in $\VVV_{\SSS}$ and $\VVV_{\SSS'''}$ respectively, then we want to show that the unitary map this teleportation scheme ultimately implements is the one that identifies these operators.
    The way we'll do this is via the induced homomorphism on the logical Pauli groups, as in \Cref{ex:teleportation_logical_pauli_group_updates}.
\end{example}

\subsection{Evolution of $\mathcal{N}(\mathcal{S})/\mathcal{S}$}\label{subsec:evolution_of_logical_pauli_group}

Let $\mathcal{S} = \langle S_1, \ldots, S_r \rangle$ be a stabilizer group, and $\mathcal{N}(\mathcal{S})/\mathcal{S} = \langle \ol{i}, \ol{\XXX}_1, \ol{\ZZZ}_1, \ldots, \ol{\XXX}_k, \ol{\ZZZ}_k \rangle $ the logical Pauli group.
In this subsection, we give a presentation for the new logical Pauli group $\mathcal{N}(\mathcal{S}')/\mathcal{S}'$ after applying a Clifford or measuring a Pauli.
We also give an induced homomorphism $\phi: \mathcal{N}(\mathcal{S}')/\mathcal{S}' \to \mathcal{N}(\mathcal{S})/\mathcal{S}$, which captures how the logical operators of the new codespace relate to logical operators of the old codespace.
Note the arrow goes the `wrong' way!
The relationship it captures is the following.
For all $L' \in \NNN(\SSS)$ and $L \in \NNN(\SSS)$ with $\overline{L} \in im(\phi)$:
\begin{equation}\label{eq:induced_homomorphism_iff_relation}
    \phi(\overline{L'}) = \overline{L} \iff (L' A_m)|_{\VVV_\SSS} = (A_m L)|_{\VVV_\SSS}
\end{equation}
where $A_m$ is the induced linear map on $\mathcal{H}$ from the last subsection.
The second equality can be equivalently phrased in terms of projectors:
\begin{equation}\label{eq:induced_homomorphism_iff_relation_via_projector}
    (L' A_m)|_{\VVV_\SSS} = (A_m L)|_{\VVV_\SSS} \iff L' A_m \Pi_{\SSS} = A_m L \Pi_{\SSS}
\end{equation}
Intuitively, this condition says that the logical action of $L'$ on the new codespace is exactly that of $L$ on the old codespace.
In fact, this induced homomorphism actually defines a logical Clifford unitary on our codespace.
For tidiness' sake we defer discussion of this and the proofs that condition~\eqref{eq:induced_homomorphism_iff_relation} is satisfied in each case to its own subsection below.

After applying Clifford unitary $U$, we've seen (\Cref{thm:stabilizer_group_update_clifford}) that the new stabilizer group is $\mathcal{S}' = U\mathcal{S}U^\dagger = \langle U S_1 U^\dagger, \ldots, U S_r U^\dagger \rangle$.
The new logical Pauli group is updated similarly.

\begin{theorem}\label{thm:logical_pauli_group_update_clifford}
    \cite[Procedure 6.3]{gottesman2024surviving} If a system with stabilizer group $\mathcal{S} = \langle S_1, \ldots, S_r \rangle$ and logical Pauli group $\mathcal{N}(\mathcal{S})/\mathcal{S} = \langle \overline{i}, \overline{\mathcal{X}}_1, \overline{\mathcal{Z}}_1, \ldots, \overline{\mathcal{X}}_k, \overline{\mathcal{Z}}_k \rangle$ undergoes evolution by a Clifford unitary $U$, the new logical Pauli group is
    \begin{equation}\label{eq:new_logical_Pauli_group}
        \mathcal{N}(\mathcal{S}')/\mathcal{S}' = U (\mathcal{N}(\mathcal{S})/\mathcal{S}) U^\dagger = \langle \ol{i}, \ol{U \XXX_1 U^\dagger}, \ol{U \ZZZ_1 U^\dagger}, \ldots, \ol{U \XXX_k U^\dagger},  \ol{U \ZZZ_k U^\dagger} \rangle
    \end{equation}

    The induced group homomorphism $\phi : \mathcal{N}(\mathcal{S}')/\mathcal{S}' \to \mathcal{N}(\mathcal{S})/\mathcal{S}$ is defined to be $\phi(\overline{P'}) = \overline{U^\dagger P' U}$.
    In this case, this is in fact an isomorphism.
    Equivalently, this can be defined by its action on generators:
    \begin{equation}
    %    \begin{aligned}
    %        \ol{i} &\mapsto \ol{i} \\
    %        \ol{\XXX}_j &\mapsto \ol{U^\dagger \XXX_j U} \text { for all $1 \leq j \leq k$} \\
    %        \ol{\ZZZ}_j &\mapsto \ol{U^\dagger \ZZZ_j U} \text { for all $1 \leq j \leq k$} \\
    %    \end{aligned}
        \ol{i} \mapsto \ol{i}, \quad
        \ol{\XXX}_j \mapsto \ol{U^\dagger \XXX_j U}, \quad
        \ol{\ZZZ}_j \mapsto \ol{U^\dagger \ZZZ_j U}, \quad
        \text { for all $1 \leq j \leq k$}
    \end{equation}
\end{theorem}

The effect of measuring a Pauli operator $P$ on the logical Pauli group $\mathcal{N}(\mathcal{S})/\mathcal{S}$ depends on which of the three cases we are in.
We first need a couple of lemmas about presentations of $\mathcal{N}(\mathcal{S})/\mathcal{S}$.

\begin{lemma}\label{lem:logical_presentation_WLOG_case_2}
    Suppose we have stabilizer group $\mathcal{S} = \langle S_1, \ldots, S_r \rangle \leq \PPP_n$ and logical Pauli group $\mathcal{N}(\mathcal{S})/\mathcal{S} = \langle \overline{i}, \overline{\mathcal{X}}_1, \overline{\mathcal{Z}}_1, \ldots, \overline{\mathcal{X}}_k, \overline{\mathcal{Z}}_k \rangle$ satisfying the canonical commutation relations.
    Let $P \in \PPP_n$ be some Hermitian Pauli that anti-commutes with at least one stabilizer $S \in \SSS$.
    We can always find a presentation of $\mathcal{N}(\mathcal{S})/\mathcal{S}$ such that $P$ commutes with all representatives $\XXX_j$ and $\ZZZ_j$, and the canonical commutation relations are preserved.
    \begin{proof}
        We know we can assume without loss of generality that $P$ anti-commutes with only $S_1$ and commutes with the remaining stabilizer generators.
        So now if $P$ anti-commutes with a representative $\XXX_j$ of $\XXX_j \SSS$, we can pull out $S_1$ from $\SSS$ to get a new representative $\XXX_j S_1$ that does commute with $P$.
        Same goes for $\ZZZ_j$.
    \end{proof}
\end{lemma}

\begin{lemma}\label{lem:logical_presentation_WLOG_case_3}
    Suppose we have stabilizer group $\mathcal{S} = \langle S_1, \ldots, S_r \rangle \leq \PPP_n$ and logical Pauli group $\mathcal{N}(\mathcal{S})/\mathcal{S} = \langle \overline{i}, \overline{\mathcal{X}}_1, \overline{\mathcal{Z}}_1, \ldots, \overline{\mathcal{X}}_k, \overline{\mathcal{Z}}_k \rangle$ satisfying the canonical commutation relations.
    Let $P \in \PPP_n$ be some Hermitian Pauli that commutes with all of $\mathcal{S}$, but $\pm P \notin \mathcal{S}$.
    That is, $P$ is in $\NNN(\SSS) - (\pm \SSS)$, so is a non-trivial element of the logical Pauli group.
    We can always find a presentation of $\mathcal{N}(\mathcal{S})/\mathcal{S}$ such that $\ol{P} = \pm \ol{\ZZZ}_j$ or $\pm \ol{\XXX}_j$ for any $1 \leq j \leq k$, and the canonical commutation relations are preserved.
    \begin{proof}
        This proof is a bit wordy, so consider looking at it in action in \Cref{ex:bacon_shor_logical_pauli_group_updates} for more clarity.
        Since $P \in \NNN(\SSS)$, we can write it as $P = i^c \XXX_1^{a_1}\ZZZ_1^{b_1} \ldots \XXX_k^{a_1}\ZZZ_k^{b_1} S$, for some $S \in \SSS$.
        And we can determine the values of the $a_j$ and $b_j$ terms via commutation relations.
        Namely, $a_j = 1$ iff $P$ and $\ZZZ_j$ anti-commute, and $b_j = 1$ iff $P$ and $\XXX_j$ anti-commute.
        So now we can just start multiplying generators of $\NNN(\SSS)/\SSS$ together until $\pm \ol{P}$ is a generator, while maintaining the canonical commutation relations between them.
        We can achieve this as follows.
        Let $a_{j_1}$ be the first $a_j$ term that equals 1.
        We will multiply $\ol{\XXX}_{j_1}$ with other generators such that at the end this equals $\pm \ol{P}$.
        First, multiply it by $\ol{\XXX}_{j_k}$ for each ${a_{j_k}}$ that equals 1.
        But to maintain the canonical commutation relations, we also multiply each $\ol{\ZZZ}_{j_k}$ by $\ol{\XXX}_{j_1}$.
        At this point, we've transformed $\ol{\XXX}_{j_1} \mapsto \ol{\XXX}_{j_1} \ldots \ol{\XXX}_{j_r}$.
        Repeat for the $\ZZZ_l$ terms; multiply $\ol{\XXX}_{j_1} \ldots \ol{\XXX}_{j_r}$ by $\ol{\ZZZ}_{l_m}$ for each $b_{l_m} = 1$, and in turn multiply each $\ol{\XXX}_{l_m}$ by $\ol{\ZZZ}_{j_1}$ to maintain the commutativity relations.
        So we're essentially done -- just note that we can now freely reorder anti-commuting pairs of generators, swap the roles of $\ol{\XXX}_j$ and $\ol{\ZZZ}_j$ within a pair, and multiply them by $\pm 1$, until $\ol{P}$ equals a specific $\pm \ol{\XXX}_j$ or $\pm \ol{\ZZZ}_j$.
    \end{proof}
\end{lemma}

Now we state the main theorem:

\begin{theorem}\label{thm:logical_pauli_group_update_measurement}
    \cite[Theorem 6.6]{gottesman2024surviving} If a system with stabilizer group $\mathcal{S} = \langle S_1, \ldots, S_r \rangle$ and logical Pauli group $\mathcal{N}(\mathcal{S})/\mathcal{S} = \langle \overline{i}, \overline{\mathcal{X}}_1, \overline{\mathcal{Z}}_1, \ldots, \overline{\mathcal{X}}_k, \overline{\mathcal{Z}}_k \rangle$ undergoes evolution by a Pauli measurement $P$, the effect is as follows.

    \textbf{Case 1:}\label{thm:logical_pauli_group_update_measurement:case1} $\pm P \in \mathcal{S}$.
    The measurement outcome is deterministic and the state is unchanged.
    The stabilizer group and logical Pauli group remain the same: $\mathcal{S}' = \mathcal{S}$ and $\mathcal{N}(\mathcal{S}')/\mathcal{S}' = \mathcal{N}(\mathcal{S})/\mathcal{S}$.
    The induced homomorphism is the identity $\phi = id$; this is trivially an isomorphism.


    \textbf{Case 2:}\label{thm:logical_pauli_group_update_measurement:case2} $P$ anticommutes with at least one element of $\mathcal{S}$.
    The measurement outcome $m \in \{+1, -1\}$ is uniformly random.
    Without loss of generality, assume $P$ anticommutes only with $S_1$ and commutes with $S_2, \ldots, S_r$.
    Also without loss of generality, assume $P$ commutes with all representatives $\XXX_j$ and $\ZZZ_j$ (\Cref{lem:logical_presentation_WLOG_case_2}).
    Then the new logical Pauli group can be presented by the same representatives:
    That is:
    \begin{equation}
        \mathcal{N}(\mathcal{S}')/\mathcal{S}' = \langle \ol{i}, \ol{\XXX}_1, \ol{\ZZZ}_1, \ldots, \ol{\XXX}_k, \ol{\ZZZ}_k \rangle
    \end{equation}
    The induced homomorphism is defined as $\phi(\overline{L'}) = \overline{L'}$, assuming without loss of generality that the representative $L'$ commutes with $S_1$ and hence all of $\SSS$.
    This is again an isomorphism.
    Equivalently, this can be defined by its action on generators:
    \begin{equation}
    %    \begin{aligned}
    %        \ol{i} &\mapsto \ol{i} \\
    %        \ol{\XXX}_j &\mapsto \ol{\XXX}_j \text { for all $1 \leq j \leq k$} \\
    %        \ol{\ZZZ}_j &\mapsto \ol{\ZZZ}_j \text { for all $1 \leq j \leq k$} \\
    %    \end{aligned}
        \ol{i} \mapsto \ol{i}, \quad
        \ol{\XXX}_j \mapsto \ol{\XXX}_j, \quad
        \ol{\ZZZ}_j \mapsto \ol{\ZZZ}_j, \quad
        \text { for all $1 \leq j \leq k$}
    \end{equation}

    \textbf{Case 3:}\label{thm:logical_pauli_group_update_measurement:case3} $P$ commutes with all of $\mathcal{S}$, but $\pm P \notin \mathcal{S}$.
    In this case, $P$ is in $\NNN(\SSS) - (\pm \SSS)$, so is a non-trivial element of the logical Pauli group $\mathcal{N}(\mathcal{S})/\mathcal{S}$.
    Without loss of generality, we can in fact assume $P = \pm \ZZZ_j$ or $\pm \XXX_j$ for some $1 \leq j \leq k$ (\Cref{lem:logical_presentation_WLOG_case_3}).
    So we're measuring one logical qubit in a Pauli basis, effectively removing a logical qubit from the system.
    Accordingly, the new logical Pauli group can be presented as
    \begin{equation}
        \mathcal{N}(\mathcal{S}')/\mathcal{S}' = \langle \ol{i}, \ol{\XXX}_1, \ol{\ZZZ}_1, \ldots, \ol{\XXX}_{j-1}, \ol{\ZZZ}_{j-1}, \ol{\XXX}_{j+1}, \ol{\ZZZ}_{j+1}, \ldots, \ol{\XXX}_{k}, \ol{\ZZZ}_{k} \rangle
    \end{equation}
    The induced homomorphism is defined as $\phi(\overline{L'}) = \overline{L'}$, assuming without loss of generality that the representative $L'$ commutes with $\mathcal{X}_j$ if $P = \pm \ZZZ_j$, or with $\ZZZ_j$ if $P = \pm \XXX_j$.
    Equivalently, it can be defined by its action on generators:
    \begin{equation}\label{eq:case_3_phi_definition_generators}
        \ol{i} \mapsto \ol{i}, \quad
        \ol{\XXX}_l \mapsto \ol{\XXX}_l, \quad
        \ol{\ZZZ}_l \mapsto \ol{\ZZZ}_l, \quad
        \text { for all $0 \leq l \leq k$ with $l \neq j$}
    \end{equation}
\end{theorem}

Note that in \hyperref[thm:logical_pauli_group_update_measurement:case2]{Case 2}, the new logical Pauli group might look identical to the old one, but it isn't!
The cosets $\overline{L}$ in the two groups are different; in the old group $\overline{L} = L\SSS$ and in the new group $\overline{L} = L\SSS'$, where $\SSS \neq \SSS'$.
Similarly, despite appearances, the induced homomorphism isn't the identity.
\hyperref[thm:logical_pauli_group_update_measurement:case3]{Case 3} is the only one in which the induced homomorphism isn't an isomorphism, since the new logical Pauli group is smaller than the original.
In both \hyperref[thm:logical_pauli_group_update_measurement:case2]{Cases 2} and \hyperref[thm:logical_pauli_group_update_measurement:case3]{3}, we said that without loss of generality we could make assumptions about the commutativity relations between $P$ and representatives $\XXX_j$ and $\ZZZ_j$, which we now spell out.

\begin{example}\label{ex:bacon_shor_logical_pauli_group_updates}
    Recall \Cref{ex:bacon_shor_stabilizer_group_updates,ex:bacon_shor_codespace_updates}, in which we started with a trivial stabilizer group on four qubits and repeatedly measured $X_1 X_2$, $X_3 X_4$, $Z_1 Z_3$ and $Z_2 Z_4$, in that order.
    We spelled out how the stabilizer group and codespace changes throughout.
    Now we'll do the same for the logical Pauli group.
    Initially, since $\SSS_{-1} = \{I\}$, the following presentation obeys the Pauli commutation relations:
    \begin{equation}
        \NNN(\SSS_{-1})/\SSS_{-1} = \langle
            \overline{i},
            \overline{X_1}, \overline{Z_1},
            \overline{X_2}, \overline{Z_2},
            \overline{X_3}, \overline{Z_3},
            \overline{X_4}, \overline{Z_4},
        \rangle
    \end{equation}
    The first measurement $X_1 X_2$ is in \hyperref[thm:stabilizer_group_update_measurement:case3]{Case 3}.
    To apply \Cref{thm:logical_pauli_group_update_measurement}, we need a new presentation for $\NNN(\SSS_{-1})/\SSS_{-1}$ in which $X_1 X_2$ is a representative $\mathcal{X}_j$ or $\mathcal{Z}_j$, and in which the Pauli commutativity relations are still obeyed.
    Following \Cref{lem:logical_presentation_WLOG_case_3}, we can write
    \begin{equation}
        \NNN(\SSS_{-1})/\SSS_{-1} = \langle
        \overline{i},
        \overline{X_1 X_2}, \overline{Z_1},
        \overline{X_2}, \overline{Z_1 Z_2},
        \overline{X_3}, \overline{Z_3},
        \overline{X_4}, \overline{Z_4},
        \rangle
    \end{equation}
    After measurement, the new group is:
    \begin{equation}
        \NNN(\SSS_{0})/\SSS_{0} = \langle
        \overline{i},
        \overline{X_2}, \overline{Z_1 Z_2},
        \overline{X_3}, \overline{Z_3},
        \overline{X_4}, \overline{Z_4},
        \rangle
    \end{equation}
    The induced homomorphism $\phi_0: \NNN(\SSS_{0})/\SSS_{0} \to \NNN(\SSS_{-1})/\SSS_{-1}$ is $\phi_0(\overline{L}) = \overline{L}$, assuming without loss of generality that $L$ commutes with $Z_1$.
    Next, we measure $X_3 X_4$, which is in \hyperref[thm:stabilizer_group_update_measurement:case3]{Case 3}.
    We again first require a new presentation:
    \begin{equation}
        \NNN(\SSS_{0})/\SSS_{0} = \langle
            \overline{i},
            \overline{X_2}, \overline{Z_1 Z_2},
            \overline{X_3 X_4}, \overline{Z_3},
            \overline{X_4}, \overline{Z_3 Z_4},
        \rangle
    \end{equation}
    Now we can apply \Cref{thm:logical_pauli_group_update_measurement}; the new logical Pauli group is
    \begin{equation}
        \NNN(\SSS_{0}')/\SSS_{0}' = \langle
            \overline{i},
            \overline{X_2}, \overline{Z_1 Z_2},
            \overline{X_4}, \overline{Z_3 Z_4},
        \rangle
    \end{equation}
    The induced homomorphism $\phi_0': \NNN(\SSS_{0}')/\SSS_{0}' \to \NNN(\SSS_{0})/\SSS_{0}$ is $\phi_0'(\overline{L'}) = \overline{L'}$, assuming without loss of generality that $L'$ commutes with $Z_3$.
    Next, we measure $Z_1 Z_3$, which is our first \hyperref[thm:stabilizer_group_update_measurement:case2]{Case 2} measurement.
    This already commutes with all coset representatives in our presentation of the logical Pauli group, so we can directly apply \Cref{thm:logical_pauli_group_update_measurement}, which says that the new group can be presented using exactly the same coset representatives as before:
    \begin{equation}
        \NNN(\SSS_{0}'')/\SSS_{0}'' = \langle
        \overline{i},
        \overline{X_2}, \overline{Z_1 Z_2},
        \overline{X_4}, \overline{Z_3 Z_4},
        \rangle
    \end{equation}
    The induced homomorphism $\phi_0'': \NNN(\SSS_{0}'')/\SSS_{0}'' \to \NNN(\SSS_{0}')/\SSS_{0}'$ is $\phi_0''(\overline{L''}) = \overline{L''}$, assuming without loss of generality that $L''$ commutes with all of $\SSS_{0}' = \langle a_0 X_1 X_2, b_0 X_1 X_3 \rangle$.
    The final measurement in round $0$ is $Z_2 Z_4$, which is in \hyperref[thm:stabilizer_group_update_measurement:case3]{Case 3}.
    To apply \Cref{thm:logical_pauli_group_update_measurement}, we need $Z_2 Z_4$ to be a representative of a coset in $\NNN(\SSS_{0}'')/\SSS_{0}''$.
    So first apply \Cref{lem:logical_presentation_WLOG_case_3} to get:
    \begin{equation}
        \NNN(\SSS_{0}'')/\SSS_{0}'' = \langle
        \overline{i},
        \overline{X_2}, \overline{Z_1 Z_2 Z_3 Z_4},
        \overline{X_2 X_4}, \overline{Z_3 Z_4},
        \rangle
    \end{equation}
    Now recall that at this point $c_0 Z_1 Z_3$ is a stabilizer, so we can `pull it out' of $\overline{Z_1 Z_2 Z_3 Z_4}$ to get a new representative $c_0 Z_2 Z_4$ of this coset.
    Now we can apply \Cref{thm:logical_pauli_group_update_measurement} to see the new group is:
    \begin{equation}
        \NNN(\SSS_{0}''')/\SSS_{0}''' = \langle
        \overline{i},
        \overline{X_2 X_4}, \overline{Z_3 Z_4},
        \rangle
    \end{equation}
    From now on, we can always present the logical Pauli group with these representatives!
    This is because future measurements are always either in \hyperref[thm:stabilizer_group_update_measurement:case1]{Case 1} or \hyperref[thm:stabilizer_group_update_measurement:case2]{Case 2}, and the representatives above commute with all of $X_1 X_2$, $X_3 X_4$, $Z_1 Z_3$ and $Z_2 Z_4$.
\end{example}

\begin{example}\label{ex:teleportation_logical_pauli_group_updates}
    Recall the teleportation scheme we looked at in \Cref{ex:teleportation_stabilizer_group_updates,ex:teleportation_codespace_updates}.
    We now look at this scheme's effect on the logical Pauli group.
    We start with $\SSS = \langle X_2 X_3, Z_2 Z_3 \rangle$, hence:
    \begin{equation}
        \NNN(\SSS)/\SSS = \langle
            \overline{i},
            \overline{X_1}, \overline{Z_1},
        \rangle
    \end{equation}
    We measure $X_1 X_2$ and get outcome $a$.
    This is a \hyperref[thm:stabilizer_group_update_measurement:case2]{Case 2} measurement, so to apply \Cref{thm:logical_pauli_group_update_measurement} we need a new presentation whose representatives commute with $X_1 X_2$.
    Noting that $Z_2 Z_3$ is a stabilizer, we can use:
    \begin{equation}
        \NNN(\SSS)/\SSS = \langle
            \overline{i},
            \overline{X_1}, \overline{Z_1 Z_2 Z_3},
        \rangle
    \end{equation}
    Now \Cref{thm:logical_pauli_group_update_measurement} says the new group can be presented by these same representatives, so:
    \begin{equation}
        \NNN(\SSS')/\SSS' = \langle
        \overline{i},
        \overline{X_1}, \overline{Z_1 Z_2 Z_3},
        \rangle
    \end{equation}
    The induced homomorphism $\phi$ maps $\phi(\overline{L'}) = \overline{L'}$, assuming without loss of generality that $L'$ commutes with all of $\SSS = \langle X_2 X_3, Z_2 Z_3 \rangle$.
    Next we measure $Z_1 Z_2$ and are in a symmetric situation; this is in \hyperref[thm:stabilizer_group_update_measurement:case2]{Case 2}, so we get random outcome $b$, and since $X_2 X_3$ is a stabilizer a suitable new presentation for the logical Pauli group is:
    \begin{equation}
        \NNN(\SSS')/\SSS' = \langle
        \overline{i},
        \overline{X_1 X_2 X_3}, \overline{Z_1 Z_2 Z_3},
        \rangle
    \end{equation}
    Hence these same representatives can be used to present the new group:
    \begin{equation}
        \NNN(\SSS'')/\SSS'' = \langle
        \overline{i},
        \overline{X_1 X_2 X_3}, \overline{Z_1 Z_2 Z_3},
        \rangle
    \end{equation}
    The induced homomorphism $\phi''$ maps $\phi(\overline{L''}) = \overline{L''}$, assuming without loss of generality that $L''$ commutes with all of $\SSS' = \langle X_2 X_3, a X_1 X_2 \rangle$.
    Since at this point $\SSS'' = \langle a X_1 X_2, b Z_1 Z_2 \rangle$, we can pull out new representatives for each coset:
    \begin{equation}
        \NNN(\SSS'')/\SSS'' = \langle
        \overline{i},
        \overline{a X_3}, \overline{b Z_3},
        \rangle
    \end{equation}
    Finally, we apply Pauli correction $U = Z_3^{\frac{1 - a}{2}} X_3^{\frac{1 - b}{2}}$.
    This gives the new group:
    \begin{equation}
        \NNN(\SSS''')/\SSS''' = \langle
        \overline{i},
        \overline{X_3}, \overline{Z_3},
        \rangle
    \end{equation}
    The induced homomorphism $\phi'''$ maps $\phi(\overline{L'''}) = \overline{L'''}$, assuming without loss of generality that $L'''$ commutes with all of $\SSS'' = \langle a X_1 X_2, b Z_1 Z_2 \rangle$.
    We can now almost see that this teleportation scheme works as intended!
    We can at least check the images of $\overline{X_3}$ and $\overline{Z_3}$ under the composition $\phi \circ \phi' \circ \phi'' \circ \phi''' \circ \phi''''$ of all the induced homomorphisms are $\overline{X_1}$ and $\overline{Z_1}$ respectively.
    That is, writing cosets in full as $L\SSS$ rather than $\overline{L}$, we have:
%    \begin{equation}
%        \begin{aligned}
%            X_3\SSS'''
%            &\xmapsto{\phi'''} a X_3\SSS'' = X_1 X_2 X_3 \SSS'' \\
%            &\xmapsto{\phi''} X_1 X_2 X_3 \SSS' = X_1 \SSS' \\
%            &\xmapsto{\phi'} X_1\SSS
%        \end{aligned}
%    \end{equation}
%    \begin{equation}
%        \begin{aligned}
%            Z_3\SSS'''
%            &\xmapsto{\phi'''} b Z_3\SSS'' = Z_1 Z_2 Z_3 \SSS'' \\
%            &\xmapsto{\phi''} Z_1 Z_2 Z_3 \SSS' = Z_1 \SSS' \\
%            &\xmapsto{\phi'} Z_1\SSS
%        \end{aligned}
%    \end{equation}
    \begin{align}
        X_3\SSS'''
        &\xmapsto{\phi'''} a X_3\SSS'' = X_1 X_2 X_3 \SSS''
        \xmapsto{\phi''} X_1 X_2 X_3 \SSS' = X_1 \SSS'
        \xmapsto{\phi'} X_1\SSS \\
        Z_3\SSS'''
        &\xmapsto{\phi'''} b Z_3\SSS'' = Z_1 Z_2 Z_3 \SSS''
        \xmapsto{\phi''} Z_1 Z_2 Z_3 \SSS' = Z_1 \SSS'
        \xmapsto{\phi'} Z_1\SSS \\
    \end{align}
    Combining this with the defining relation \Cref{eq:induced_homomorphism_iff_relation} for the induced homomorphism, this says intuitively that $\overline{X_3}$ and $\overline{Z_3}$ at the end of the teleportation scheme perform the same logical operation as $\overline{X_1}$ and $\overline{Z_1}$ at the start of the scheme.
    Rigorously, we'll show in the next subsection that whenever the induced homomorphism is in fact an isomorphism, it uniquely determines a logical Clifford unitary action on our codespace.
    In this case here, this logical Clifford is the identity, so the state of qubit 1 really is being teleported to qubit 3.
\end{example}

\subsection{The induced homomorphism $\phi$}\label{subsec:the-induced-homomorphism}

As promised, we now show two things.
First, we'll show that whenever this homomorphism is in fact an isomorphism then it uniquely determines a logical Clifford unitary that's applied to our codespace.
Then, we show that in each of the four settings above (applying a Clifford and the three cases of Pauli measurements), the induced homomorphism $\phi: \NNN(\SSS') / \SSS' \to \NNN(\SSS) / \SSS$ satisfies the required iff relationship~\eqref{eq:induced_homomorphism_iff_relation}.

So, first assume the homomorphism is an isomorphism, which occurs iff all operations are Clifford unitaries, or \hyperref[thm:stabilizer_group_update_measurement:case1]{Case 1} or \hyperref[thm:stabilizer_group_update_measurement:case2]{2} Pauli measurements.
Recall the defining iff relation for the induced homomorphism as in \Cref{eq:induced_homomorphism_iff_relation}.
Namely, for all $L' \in \NNN(\SSS)$ and $L \in \NNN(\SSS)$:
\begin{equation}
    \phi(\overline{L'}) = \overline{L}
    \iff (L' A_m)|_{\VVV_\SSS} = (A_m L)|_{\VVV_\SSS}
    \iff L' A_m \Pi_{\SSS} = A_m L \Pi_{\SSS}
\end{equation}
(We've omitted the extra condition that $\overline{L} \in im(\phi)$ because it's irrelevant if $\phi$ is an isomorphism).
In the cases we're considering, we've seen that the linear map $A_m|_{\VVV_\SSS} = A_m \Pi_\SSS$ applied to the codespace is unitary.
Also note that any $\overline{L} \in \NNN(\SSS)/\SSS$ by definition commutes with all $S \in \SSS$, so in particular commutes with $\Pi_\SSS$.
Therefore $A_m L \Pi_{\SSS} = A_m \Pi_{\SSS} L$, and we can multiply both sides of $L' A_m \Pi_{\SSS} = A_m L \Pi_{\SSS} = A_m \Pi_{\SSS} L$ by the adjoint of the unitary $A_m \Pi_\SSS$ to get $L' = (A_m \Pi_{\SSS}) L (A_m \Pi_{\SSS})^\dagger$.
So we're saying that the net effect of applying this physical Clifford unitary or performing this physical Pauli measurement is that a unitary map $A_m \Pi_{\SSS}$ is applied to our codespace, which acts via conjugation on logical Pauli operators.
Furthermore, it conjugates these Paulis $L$ back to Paulis $L' = (A_m \Pi_{\SSS}) L (A_m \Pi_{\SSS})^\dagger$.
It is then a standard result (e.g.\ \Ccite[Theorem 6.1]{gottesman2024surviving}) that this map $A_m \Pi_{\SSS}$ must be a uniquely determined logical Clifford (unique up to global phase, at least).

\begin{example}\label{ex:teleportation_logical_clifford}
    We can finally finish our running example of a teleportation scheme (\Cref{ex:teleportation_logical_pauli_group_updates}).
    Our initial logical Pauli group was $\pres{\ol{i}, \ol{\XXX} = \ol{X_1}, \ol{\ZZZ} = \ol{Z_1}}$, and our final logical Pauli group is $\pres{\ol{i}, \ol{\XXX} = \ol{X_3}, \ol{\ZZZ} = \ol{Z_3}}$.
    We showed that the overall induced isomorphism sends:
    \begin{equation}
        \XXX\SSS''' = X_3\SSS''' \mapsto X_1\SSS = \XXX\SSS
        \quad\text{and}\quad
        \ZZZ\SSS''' = Z_3\SSS''' \mapsto Z_1\SSS = \ZZZ\SSS
    \end{equation}
    Taking the inverse of this, we're saying that overall a logical unitary $U$ is applied that conjugates $\ol{\XXX} \mapsto U\ol{\XXX} U^\dagger = \ol{\XXX}$ and $\ol{\ZZZ} \mapsto U\ol{\ZZZ} U^\dagger = \ol{\ZZZ}$.
    One Clifford that has this action is the identity map, and since $U$ is uniquely determined (up to global phase), $U$ must in fact be the identity.
\end{example}

Finally, we give the proofs that the induced homomorphism satisfies the defining iff relation \Cref{eq:induced_homomorphism_iff_relation} in each case.
Throughout, we'll make use of the following lemma.

\begin{lemma}\label{lem:paulis_equal_on_codespace_are_equal_up_to_stabilizer}
    Given stabilizer group $\SSS$, if Paulis $P$ and $Q$ satisfy $P|_{\VVV_\SSS} = Q|_{\VVV_{\SSS}}$, i.e. $P \Pi_{\SSS} = Q \Pi_\SSS$, then we must have $P\SSS = Q\SSS$, i.e. $P = QS$ for some $S \in \SSS$.
    \begin{proof}
        Define $R = Q^\dagger P$.
        Since $P \Pi_{\SSS} = Q \Pi_\SSS$ we have:
        \begin{equation}
            R \Pi_\SSS = Q^\dagger P \Pi_{\SSS} = Q^\dagger Q \Pi_\SSS = \Pi_\SSS
        \end{equation}
        That is, $R$ is a stabilizer for all $\ket{\psi} \in \VVV_\SSS$, in that $R\ket{\psi} = \ket{\psi}$.
        A defining property of a non-trivial Pauli stabilizer code is that any Pauli that stabilizes all codespace elements must be in $\SSS$ (e.g.\ \Ccite[Proposition 3.2]{gottesman2024surviving}).
        So $R \in \SSS$, hence $P = QQ^\dagger P = QR$, as required.
    \end{proof}
\end{lemma}

We start with the Clifford unitary case.

\begin{proposition}
    After applying a Clifford $U$, the induced homomorphism $\phi$ satisfies \Cref{eq:induced_homomorphism_iff_relation}.
    That is, $\phi(\overline{L'}) = \overline{L} \iff (L' A_m)|_{\VVV_\SSS} = (A_m L)|_{\VVV_\SSS}$, where $A_m = U$.
    \begin{proof}
        In one direction, assume $\overline{L} = \phi(\overline{L'}) = \overline{U^\dagger L' U}$.
        Hence $L = U^\dagger L' U S$ for some $S \in \SSS$.
        So the right hand side of the iff condition becomes $(L' U)|_{\VVV_\SSS} = (U U^\dagger L' U S)|_{\VVV_\SSS}$, and these are indeed equal!
        Conversely, suppose $L$ and $L'$ satisfy the right hand side $(L' U)|_{\VVV_\SSS} = (U L)|_{\VVV_\SSS}$.
        Therefore $(U^\dagger L' U)|_{\VVV_\SSS} = (U L)|_{\VVV_\SSS}$,
        and hence $U^\dagger L' U \SSS = P \SSS$ by \Cref{lem:paulis_equal_on_codespace_are_equal_up_to_stabilizer}.
        So indeed $\phi(\overline{L'}) = \overline{U^\dagger L' U} = \overline{L}$, as required.
    \end{proof}
\end{proposition}

Now for the Pauli measurements.
Hopefully it's clear that in \hyperref[thm:stabilizer_group_update_measurement:case1]{Case 1} there's really nothing to prove!
But \hyperref[thm:stabilizer_group_update_measurement:case2]{Case 2} and \hyperref[thm:stabilizer_group_update_measurement:case3]{Case 3} are non-trivial.

\begin{proposition}
    In \hyperref[thm:stabilizer_group_update_measurement:case2]{Case 2}, the induced homomorphism $\phi$ satisfies \Cref{eq:induced_homomorphism_iff_relation}.
    That is, after measuring $P$ and getting outcome $m$, we have $\phi(\overline{L'}) = \overline{L} \iff (L' A_m)|_{\VVV_\SSS} = (A_m L)|_{\VVV_\SSS}$, where $A_m = \Pi_{mP}$.
    \begin{proof}
        In one direction, assume $\overline{L} = \phi(\overline{L'}) \coloneqq \overline{L'}$, for $L \in \NNN(\SSS)$ and $L' \in \NNN(\SSS')$.
        Hence $L = L'S$ for some $S \in \SSS$.
        The right hand side of the iff condition becomes $(L' \Pi_{mP})|_{\VVV_\SSS} = (\Pi_{mP} L)|_{\VVV_\SSS}$.
        Subbing $L = L'S$ into the right hand side and absorbing the $S$, this becomes $(L' \Pi_{mP})|_{\VVV_\SSS} = (\Pi_{mP} L')|_{\VVV_\SSS}$.
        Now note that since $L' \in \NNN(\SSS')$ and $mP \in \SSS'$, we know $L'$ and $P$ commute, hence so do $L'$ and $\Pi_{mP}$, so $(L' \Pi_{mP})|_{\VVV_\SSS} = (\Pi_{mP} L')|_{\VVV_\SSS}$ as required.

        Conversely, suppose $L \in \NNN(\SSS)$ and $L' \in \NNN(\SSS')$ satisfy the right hand side equality $(L'
        \Pi_{mP})|_{\VVV_\SSS} = (\Pi_{mP} L)|_{\VVV_\SSS}$.
%        Conversely, suppose $L \in \NNN(\SSS)$ and $L' \in \NNN(\SSS')$ satisfy the right hand side equality $(L'
%        \Pi_{mP})|_{\VVV_\SSS} = (\Pi_{mP} L)|_{\VVV_\SSS}$,
%        or equivalently $\Pi_{\SSS} L' \Pi_{mP} \Pi_{\SSS} = \Pi_{\SSS} \Pi_{mP} L \Pi_{\SSS}$.
        By definition $L'$ commutes with $mP$ because $mP \in \SSS'$, hence $(\Pi_{mP} L')|_{\VVV_\SSS} = (\Pi_{mP} L)|_{\VVV_\SSS}$.
        Now, since $L'$ and $L$ restricted to $\VVV_\SSS$ have images again inside $\VVV_\SSS$, and we've shown $\Pi_{mP}|_{\VVV_\SSS}$ is a unitary up to normalisation, we can cancel the $\Pi_{mP}$ from both sides to get $L'|_{\VVV_\SSS} = L|_{\VVV_\SSS}$
        Then by \Cref{lem:paulis_equal_on_codespace_are_equal_up_to_stabilizer} we have $L'\SSS = L\SSS$, i.e. $\overline{L} = \phi(\overline{L'}) \coloneqq \overline{L'}$ as required.
    \end{proof}
\end{proposition}

\begin{proposition}
    In \hyperref[thm:stabilizer_group_update_measurement:case3]{Case 3}, the induced homomorphism $\phi$ satisfies \Cref{eq:induced_homomorphism_iff_relation}.
    That is, after measuring $P$ and getting outcome $m$, we have $\phi(\overline{L'}) = \overline{L} \iff (L' A_m)|_{\VVV_\SSS} = (A_m L)|_{\VVV_\SSS}$, where $A_m = \Pi_{mP}$.
    \begin{proof}
        The first direction is proved word for word as it was in the previous proposition.
%        The first direction goes exactly as in the previous proposition.
%        Suppose $\overline{L} = \phi(\overline{L'}) \coloneqq \overline{L'}$, for $L \in \NNN(\SSS)$ and $L' \in \NNN(\SSS')$.
%        Hence $L = L'S$ for some $S \in \SSS$.
%        The right hand side of the iff condition becomes $(L' \Pi_{mP})|_{\VVV_\SSS} = (\Pi_{mP} L)|_{\VVV_\SSS}$.
%        Subbing $L = L'S$ into the right hand side and absorbing the $S$, this becomes $(L' \Pi_{mP})|_{\VVV_\SSS} = (\Pi_{mP} L')|_{\VVV_\SSS}$.
%        Now note that since $L' \in \NNN(\SSS')$ and $mP \in \SSS'$, we know $L'$ and $P$ commute, hence so do $L'$ and $\Pi_{mP}$, so $(L' \Pi_{mP})|_{\VVV_\SSS} = (\Pi_{mP} L')|_{\VVV_\SSS}$ as required.
        Conversely, suppose $L' \in \NNN(\SSS')$ and $L \in \NNN(\SSS)$ with $\overline{L} \in im(\phi)$ satisfy the right hand side equality $(L' \Pi_{mP})|_{\VVV_\SSS} = (\Pi_{mP} L)|_{\VVV_\SSS}$, or equivalently $L' \Pi_{mP} \Pi_{\SSS} = \Pi_{mP} L \Pi_{\SSS}$.
        By \Cref{lem:if_in_im_phi_then_commutes_with_P}, since $L \in im(\phi)$ it must commute with $P$, hence with $\Pi_{mP}$, so we get $L' \Pi_{mP} \Pi_{\SSS} = L \Pi_{mP} \Pi_{\SSS}$.
        Next, note $\Pi_{mP} \Pi_{\SSS}$ is exactly $\Pi_{\SSS'}$, so we have $L' \Pi_{\SSS'} = L \Pi_{\SSS'}$, and hence by \Cref{lem:paulis_equal_on_codespace_are_equal_up_to_stabilizer} $L' = LS'$ for some $S' \in \SSS'$.
        But the only difference between $\SSS$ and $\SSS'$ is that the latter has the extra generator $mP$.
        So we can write $S' = S(mP)^a$ for $a \in \{0, 1\}$ and $S \in \SSS$.
        We want to show that $\overline{L} = \phi(\overline{L'}) \coloneqq \overline{L'}$.
        In the definition of $\phi(\overline{Q'})$, we assumed without loss of generality that the representative $Q'$ commuted with $\XXX_k$; if not, we can multiply it by the stabilizer $m\ZZZ_k = mP$.
        Hence when applying $\phi$ to the coset containing $L' = LS(mP)^a$ above, we take the representative $L' = LS$.
        Thus $\phi(\overline{L'}) = \phi(\overline{LS}) = \overline{LS} = \overline{L}$ as required.
    \end{proof}
\end{proposition}

\begin{lemma}\label{lem:if_in_im_phi_then_commutes_with_P}
    In \hyperref[thm:stabilizer_group_update_measurement:case3]{Case 3}, if $\overline{L} \in im(\phi)$ then $L$ commutes with $P$.
    \begin{proof}
        Without loss of generality we assume $P = \ZZZ_k$.
        $L \in \NNN(\SSS)$ can be written $i^c \XXX_1^{a_1} \ZZZ_1^{b_1} \ldots \XXX_k^{a_k} \ZZZ_k^{b_k} S$, where $S \in \SSS$.
        This anti-commutes with $P = \ZZZ_k$ iff $a_k = 1$.
        But no such $L$ is in $im(\phi)$, because $im(\phi) = \langle \ol{i}, \ol{\XXX}_1, \ol{\ZZZ}_1, \ldots, \ol{\XXX}_{k-1}, \ol{\ZZZ}_{k-1} \rangle \leq \NNN(\SSS)/\SSS$, as can be seen from the definition of $\phi$ in terms of generators in \Cref{eq:case_3_phi_definition_generators}.
    \end{proof}
\end{lemma}

